% Section 4.3, from lectures/L04/README.md: what you should be able to do, the questions to test
% yourself with, and the next lecture.

\section{Sammanfattning}
\label{c4:sec:summary}

\needspace{6\baselineskip}
\subsection*{Det här ska du kunna nu}
Efter det här kapitlet ska du kunna:
\begin{itemize}
  \item Rita ett Karnaughdiagram för en funktion med två, tre eller fyra ingångar, med rubrikerna i
    Graykod.
  \item Ringa in grupper av ettor efter reglerna, också över kanter och i hörn, och läsa av ett
    minimerat uttryck på summa av produkter-form.
  \item Använda don't care för kombinationer som aldrig förekommer.
  \item Syntetisera ett kombinatoriskt nät från en beskrivning i ord: sanningstabell,
    Karnaughdiagram, minimerat uttryck, grindnät, kontroll.
  \item Göra om ett nät till enbart NAND-grindar eller enbart NOR-grindar.
  \item Hitta ben 1, VCC, GND och varje grinds ben på en 74HC-krets, med hjälp av databladet.
  \item Räkna ut ett förkopplingsmotstånd för en lysdiod, och förklara varför en ingång aldrig får
    lämnas flytande.
  \item Planera en uppkoppling med en grindtilldelning och en kopplingslista.
\end{itemize}

\needspace{6\baselineskip}
\subsection*{Frågor att testa dig själv med}
\begin{enumerate}
  \item Varför måste rubrikerna i ett Karnaughdiagram stå i Graykod? Vad går fel med vanlig binär
    ordning?
  \item En grupp om fyra rutor i ett diagram med fyra ingångar: hur många ingångar blir kvar i
    termen?
  \item Ett diagram där ettorna bildar ett schackmönster går inte att förenkla. Vilken grind
    beräknar en sådan funktion?
  \item När får en X-ruta ingå i en grupp, och måste den täckas?
  \item Ett NAND-nät har ofta fler grindar än AND-OR-nätet för samma funktion. Varför kan det ändå
    vara billigare att bygga?
  \item Vilket ben är VCC och vilket är GND på en 74HC08? Och på en 74HC02? Vad skiljer 74HC02 från
    de andra?
  \item Varför är en okopplad ingång på en HC-krets ett fel, även om kretsen verkar fungera?
\end{enumerate}

\needspace{6\baselineskip}
\subsection*{Nästa kapitel}
Kapitel~\ref{c5:ch} hör ihop med Labb~2 (bilaga~\ref{app:lab2}):
\begin{itemize}
  \item \textbf{Labb~2-1}: grindnät med AND-, OR- och inverterarkretsar på kopplingsdäck.
  \item \textbf{Labb~2-2}: samma nät med enbart NAND och enbart NOR.
  \item Felsökning av ett grindnät som inte gör vad det ska, systematiskt i stället för på måfå.
\end{itemize}
